\chapter{\IfLanguageName{dutch}{Requirementsanalyse}{Requirements Analysis}}%
\label{ch:requirementsanalyse}

In dit hoofdstuk bespreken we modellen die we in de literatuurstudie al zijn tegengekomen en bespreken we de functionele vereisten, prestatievereisten en vereisten op het gebied van robuustheid en generaliseerbaarheid. Zo kunnen we de beste modellen voor onze toepassing selecteren.

\section{Functionele vereisten}

Op basis van de literatuurstudie kunnen we de volgende functionele vereisten opstellen. Het model moet:
\begin{itemize}
    \item producten kunnen detecteren in schapbeelden met hoge objectdensiteit.
    \item individuele producten kunnen lokaliseren door gebruik van bounding boxes.
    \item producten correct kunnen identificeren op SKU-niveau.
    \item kunnen omgaan met overlappende producten of scheve/beschadigde producten.
    \item visueel sterk gelijkende producten kunnen onderscheiden van elkaar.
    \item meerdere producten tegelijk kunnen verwerken binnen één beeld.
\end{itemize}


\section{Prestatievereisten}
\label{sec:prestatievereisten}

Op basis van prestatiedrempels kunnen we modellen objectief vergelijken en bekijken. We kiezen onze modellen op basis van volgende criteria:
\begin{itemize}
    \item Het model moet een mAP-score behalen op de testdataset van minstens 85\% en mAP50-95 van minstens 50\%.
    \item Het model moet een hoge recall behalen van minstens 80\%, aangezien het missen van producten vermeden moet worden.
    \item Het model moet een hoge precision behalen die minstens 85\% bedraagt, aangezien het model accuraat moet zijn over de herkenning van de producten.
    \item Voor classificatie moet een hoge recall@1 (Top-1 accuracy) worden nagestreefd die een drempelwaarde heeft van minstens 75\%.
    \item Recall@5 (Top-5 accuracy) wordt gebruikt als aanvullende metric voor visueel gelijkaardige producten.
\end{itemize}

\section{Robuustheid en generaliseerbaarheid}

Retailomgevingen variëren sterk, daarom moet het model robuust zijn tegen variaties in beeldcondities. Hiervoor kijken we naar:
\begin{itemize}
    \item variaties in belichting,
    \item verschillende camerahoeken en afstanden,
    \item variërende productdichtheid in schappen,
    \item en veranderingen in verpakking en productdesign.
\end{itemize}

De modellen moeten goed werken op beelden die lijken op de trainingsdata maar het is ook zeer belangrijk dat de modellen ook goed presteren op nieuwe of andere situaties.

\section{Vertaling van vereisten naar modelkeuze en pipeline}

De bovenstaande vereisten worden rechtstreeks vertaald naar de gekozen architectuur van de proof-of-concept. Geen enkel afzonderlijk model voldoet aan alle functionele en prestatievereisten tegelijk. Daarom wordt een combinatie van modellen gebruikt, waarbij elke component een specifieke taak uitvoert.

Voor objectdetectie kijken we naar YOLO-modellen. Deze modellen zijn geschikt voor real-time detectie en kunnen meerdere objecten tegelijk verwerken binnen één afbeelding. Het gebruik van bounding boxes sluit aan bij de functionele vereiste om individuele producten te lokaliseren. Er moet worden afgewogen tussen verschillende YOLO-varianten, waarbij de trade-off tussen nauwkeurigheid en snelheid in overweging wordt genomen.

Voor productidentificatie wordt niet gewerkt met een klassiek classificatiemodel, maar met een embedding-gebaseerde aanpak. In plaats van een vaste set klassen te voorspellen, wordt elk product voorgesteld als een vector in een hoge-dimensionale ruimte. Deze aanpak maakt het mogelijk om flexibel om te gaan met grote aantallen producten en ondersteunt het onderscheiden van visueel gelijkaardige items.

De embeddings worden gegenereerd met een combinatie van CLIP en Florence-2. CLIP levert een sterke algemene visuele representatie en presteert goed op globale kenmerken zoals kleur en merkidentiteit. Florence-2 vult dit aan met meer gedetailleerde visuele informatie, waardoor subtiele verschillen beter vastgelegd worden. Dit is noodzakelijk om te voldoen aan de vereiste om producten op SKU-niveau te onderscheiden.

Daarnaast wordt OCR toegepast via Florence-2 om tekstinformatie uit productverpakkingen te extraheren. Deze tekst wordt omgezet naar een embedding met behulp van de text encoder van CLIP. Hierdoor wordt expliciete productinformatie zoals merknamen en labels geïntegreerd in de representatie. Dit versterkt de identificatie, vooral bij producten die visueel sterk op elkaar lijken.

De combinatie van visuele en tekstuele embeddings wordt samengevoegd tot één vectorrepresentatie. Deze fusie zorgt ervoor dat verschillende informatiebronnen samen bijdragen aan de uiteindelijke vergelijking tussen producten.

Voor het zoeken naar overeenkomsten wordt gebruik gemaakt van FAISS. Deze component maakt het mogelijk om snel en efficiënt de meest gelijkaardige producten te vinden in een grote dataset. In plaats van een directe classificatie wordt gewerkt met een top-k benadering, waarbij meerdere mogelijke matches worden teruggegeven. Dit sluit aan bij het gebruik van Top-5 accuracy als aanvullende metric.

De gekozen pipeline ondersteunt verwerking van meerdere producten binnen één beeld. Elke detectie wordt afzonderlijk verwerkt, waardoor schaalbaarheid behouden blijft bij hoge objectdensiteit. Tegelijk blijft de kwaliteit afhankelijk van de nauwkeurigheid van de detectiestap en de representatiekracht van de embeddings.

Op deze manier wordt een systeem opgebouwd dat de functionele vereisten afdekt en tegelijkertijd rekening houdt met de beperkingen van individuele modellen.